\documentclass[11pt]{article}
\usepackage{amsmath,amssymb,amsthm}
\usepackage{enumitem}
\usepackage{hyperref}
\usepackage{geometry}
\geometry{margin=1in}
\newtheorem{theorem}{Theorem}[section]
\newtheorem{lemma}{Lemma}[section]
\newtheorem{assumption}{Assumption}[section]

\begin{document}

\title{Trust‑Region Gradient Method \\ Detailed Description, Convergence Theory and Full Proof}
\author{}
\date{}
\maketitle

\tableofcontents
\newpage

%%%%%%%%%%%%%%%%%%%%%%%%%%%%%%%%%%%%%%%%%%%%%%%%%%%%%%%%%%%%%%%%%%%%%%%%%%%%
\section*{Algorithm Name}
\addcontentsline{toc}{section}{Algorithm Name}
\begin{center}
\textbf{Trust‑Region Gradient (TRG) Method}
\end{center}
%%%%%%%%%%%%%%%%%%%%%%%%%%%%%%%%%%%%%%%%%%%%%%%%%%%%%%%%%%%%%%%%%%%%%%%%%%%%

\section*{Mathematical Formulation}
\addcontentsline{toc}{section}{Mathematical Formulation}
Let $f:\mathbb{R}^{n}\rightarrow\mathbb{R}$ be a continuously differentiable (possibly non‑convex) objective.  
At iteration $k$ we have a current iterate $x_{k}\in\mathbb{R}^{n}$ and a trust‑region radius $\Delta_{k}>0$.
Define the (quadratic) model
\[
m_{k}(p)\;:=\;f(x_{k})+g_{k}^{\top}p+\frac12\,p^{\top}B_{k}p,
\qquad 
g_{k}:=\nabla f(x_{k}),\;
B_{k}:=\nabla^{2}f(x_{k})\;(\text{or an approximation}).
\]

The subproblem is
\begin{equation}
\label{eq:subprob}
\min_{p\in\mathbb{R}^{n}}\; m_{k}(p)\quad\text{s.t.}\quad \|p\|\le \Delta_{k}.
\end{equation}
We solve \eqref{eq:subprob} only \emph{approximately}.  
A classical inexpensive approximation is the \emph{Cauchy point}:
\[
p_{k}^{C}\;:=\; -\frac{\Delta_{k}}{\|g_{k}\|}\,g_{k}\qquad(\text{if }g_{k}\neq0),
\]
which is the minimizer of $m_{k}$ along the steepest‑descent direction inside the ball.

Define the \emph{actual reduction}
\[
\operatorname{ared}_{k}\;:=\;f(x_{k})-f(x_{k}+p_{k}),
\]
and the \emph{predicted reduction}
\[
\operatorname{pred}_{k}\;:=\;m_{k}(0)-m_{k}(p_{k}).
\]
The ratio
\[
\rho_{k}\;:=\;\frac{\operatorname{ared}_{k}}{\operatorname{pred}_{k}}
\]
governs acceptance of the step and adaptation of $\Delta_{k}$.

\paragraph{Hyper‑parameters}
\begin{itemize}[noitemsep]
\item $\eta_{1}\in(0,1)$ – lower bound for accepting a step (typical $0.1$).
\item $\eta_{2}\in(\eta_{1},1)$ – bound for enlarging the region (typical $0.75$).
\item $\gamma_{\text{inc}}>1$ – factor for increasing $\Delta_{k}$ (e.g. $2$).
\item $\gamma_{\text{dec}}\in(0,1)$ – factor for decreasing $\Delta_{k}$ (e.g. $0.5$).
\item $\Delta_{\min}>0$, $\Delta_{\max}>0$ – allowed radius interval.
\end{itemize}

\section*{Pseudocode}
\addcontentsline{toc}{section}{Pseudocode}
\begin{enumerate}
\item \textbf{Input:} $x_{0}$, $\Delta_{0}>0$, $\eta_{1},\eta_{2},\gamma_{\text{inc}},\gamma_{\text{dec}},\Delta_{\min},\Delta_{\max}$, tolerance $\varepsilon$.
\item \textbf{For} $k=0,1,2,\dots$ until $\|g_{k}\|\le\varepsilon$ \textbf{do}
\begin{enumerate}
\item Compute $g_{k}=\nabla f(x_{k})$ (and optionally $B_{k}$).
\item Form the Cauchy point 
\[
p_{k}= -\frac{\Delta_{k}}{\|g_{k}\|}\,g_{k}\quad (\text{if }g_{k}\neq0),\qquad p_{k}=0\text{ otherwise}.
\]
\item Compute $\operatorname{ared}_{k}=f(x_{k})-f(x_{k}+p_{k})$ and 
$\operatorname{pred}_{k}=m_{k}(0)-m_{k}(p_{k})$.
\item Set $\rho_{k}= \operatorname{ared}_{k}/\operatorname{pred}_{k}$.
\item \textbf{If} $\rho_{k}\ge\eta_{1}$ \textbf{then} $x_{k+1}=x_{k}+p_{k}$ \textbf{else} $x_{k+1}=x_{k}$.
\item Update the radius:
\[
\Delta_{k+1}= \begin{cases}
\min\{\gamma_{\text{inc}}\Delta_{k},\;\Delta_{\max}\} & \text{if }\rho_{k}\ge\eta_{2},\\[4pt]
\max\{\gamma_{\text{dec}}\Delta_{k},\;\Delta_{\min}\} & \text{if }\rho_{k}<\eta_{1},\\[4pt]
\Delta_{k} & \text{otherwise.}
\end{cases}
\]
\end{enumerate}
\item \textbf{Output:} $x_{k}$.
\end{enumerate}

\section*{Dry Run Example}
\addcontentsline{toc}{section}{Dry Run Example}
Consider the scalar quadratic 
\[
f(w)=(w-3)^{2},\qquad w\in\mathbb{R}.
\]
We initialise $w_{0}=0$, set a \emph{fixed} trust‑region radius $\Delta_{k}=1$ (i.e. $\gamma_{\text{inc}}=\gamma_{\text{dec}}=1$), and use the Cauchy point.  
The gradient is $g(w)=2(w-3)$.

\begin{center}
\begin{tabular}{c|c|c|c|c}
Iteration $k$ & $w_{k}$ & $g_{k}=2(w_{k}-3)$ & $p_{k}= -\dfrac{\Delta}{|g_{k}|}g_{k}$ & $f(w_{k})$\\ \hline
0 & $0$ & $-6$ & $-\dfrac{1}{6}(-6)=1$ & $9$\\
1 & $1$ & $-4$ & $-\dfrac{1}{4}(-4)=1$ & $4$\\
2 & $2$ & $-2$ & $-\dfrac{1}{2}(-2)=1$ & $1$\\
3 & $3$ & $0$  & $0$                         & $0$
\end{tabular}
\end{center}
After three iterations the method reaches the global minimiser $w^{*}=3$ with $f(w^{*})=0$.

%%%%%%%%%%%%%%%%%%%%%%%%%%%%%%%%%%%%%%%%%%%%%%%%%%%%%%%%%%%%%%%%%%%%%%%%%%%%
\section*{Assumptions}
\addcontentsline{toc}{section}{Assumptions}
\begin{assumption}[Smoothness]
\label{as:smooth}
$f$ is continuously differentiable and its gradient is $L$‑Lipschitz, i.e.
\[
\|\nabla f(x)-\nabla f(y)\|\le L\|x-y\|\quad\forall x,y\in\mathbb{R}^{n}.
\]
\end{assumption}

\begin{assumption}[Lower Boundedness]
\label{as:lb}
There exists $f_{\inf}>-\infty$ such that $f(x)\ge f_{\inf}$ for all $x$.
\end{assumption}

\begin{assumption}[Model Accuracy]
\label{as:model}
For all $k$ and all $p$ with $\|p\|\le\Delta_{k}$,
\[
|f(x_{k}+p)-m_{k}(p)|\le \frac{L}{2}\|p\|^{2}.
\]
(Consequences of the $L$‑Lipschitz gradient.)
\end{assumption}

\begin{assumption}[Radius Bounds]
\label{as:radius}
There exist $0<\Delta_{\min}\le\Delta_{\max}<\infty$ such that
$\Delta_{k}\in[\Delta_{\min},\Delta_{\max}]$ for all $k$ (enforced by the algorithm).
\end{assumption}

%%%%%%%%%%%%%%%%%%%%%%%%%%%%%%%%%%%%%%%%%%%%%%%%%%%%%%%%%%%%%%%%%%%%%%%%%%%%
\section*{Main Convergence Theorem}
\addcontentsline{toc}{section}{Main Convergence Theorem}
\begin{theorem}
\label{thm:global}
Assume \ref{as:smooth}–\ref{as:radius} hold and that the subproblem \eqref{eq:subprob} is solved at least to the Cauchy point, i.e.
\[
m_{k}(0)-m_{k}(p_{k})\;\ge\;\frac{1}{2}\,\frac{\|g_{k}\|}{\|B_{k}\|+ \frac{\|g_{k}\|}{\Delta_{k}}}\,\|g_{k}\|\,\Delta_{k}.
\tag{Cauchy\;decrease}
\]
Then the sequence $\{x_{k}\}$ generated by the TRG method satisfies
\[
\liminf_{k\to\infty}\|\nabla f(x_{k})\|=0.
\]
Moreover, for any $\varepsilon>0$, the number of successful iterations $k$ with 
$\|\nabla f(x_{k})\|>\varepsilon$ is bounded by
\[
\mathcal{O}\!\left(\frac{L\,(f(x_{0})-f_{\inf})}{\varepsilon^{2}}\right).
\]
\end{theorem}

\section*{Theoretical Properties}
\addcontentsline{toc}{section}{Theoretical Properties}
\begin{enumerate}[noitemsep]
\item \textbf{Stability.} The radius adaptation guarantees that $\Delta_{k}$ never collapses to zero (by \ref{as:radius}) and that steps are taken only when the model predicts a sufficient decrease.
\item \textbf{Hyper‑parameter sensitivity.} Larger $\eta_{1}$ makes the algorithm more conservative (fewer accepted steps), while larger $\gamma_{\text{inc}}$ accelerates expansion once the model is trustworthy. The $O(1/\varepsilon^{2})$ bound is independent of the concrete choices as long as $0<\eta_{1}<\eta_{2}<1$ and $0<\gamma_{\text{dec}}<1<\gamma_{\text{inc}}$.
\item \textbf{Comparison.} Compared with classical (unconstrained) gradient descent with fixed stepsize $\alpha$, the trust‑region method enjoys \emph{global} convergence under merely Lipschitz gradients, whereas gradient descent may diverge for non‑convex $f$ if $\alpha$ is not carefully tuned.
\end{enumerate}

%%%%%%%%%%%%%%%%%%%%%%%%%%%%%%%%%%%%%%%%%%%%%%%%%%%%%%%%%%%%%%%%%%%%%%%%%%%%
\section*{Preliminaries (Definitions and Lemmas)}
\addcontentsline{toc}{section}{Preliminaries}
We first collect two elementary lemmas.

\begin{lemma}[Descent Lemma]\label{lem:descent}
If $\nabla f$ is $L$‑Lipschitz, then for any $x$ and $p$,
\[
f(x+p)\le f(x)+\nabla f(x)^{\top}p+\frac{L}{2}\|p\|^{2}.
\]
\end{lemma}
\begin{proof}
Standard consequence of the integral form of the gradient and the Lipschitz condition.
\end{proof}

\begin{lemma}[Cauchy decrease]\label{lem:cauchy}
Let $g\neq0$ and let $p^{C}= -\frac{\Delta}{\|g\|}g$.
For any symmetric $B$,
\[
m(0)-m(p^{C})\;\ge\;\frac{1}{2}\,\frac{\|g\|^{2}}{\|B\|+\frac{\|g\|}{\Delta}}\,\Delta .
\]
\end{lemma}
\begin{proof}
The quadratic model along the ray $p=-\alpha g$ is
\[
\phi(\alpha)= -\alpha\|g\|^{2}+\frac12\alpha^{2}g^{\top}Bg.
\]
The minimiser of $\phi$ on $[0,\Delta/\|g\|]$ is either at the unconstrained optimum
$\alpha^{\star}= \frac{\|g\|^{2}}{g^{\top}Bg}$ or at the boundary $\alpha=\Delta/\|g\|$.
A routine case analysis yields the stated lower bound.
\end{proof}

%%%%%%%%%%%%%%%%%%%%%%%%%%%%%%%%%%%%%%%%%%%%%%%%%%%%%%%%%%%%%%%%%%%%%%%%%%%%
\section*{Main Proof (Step‑by‑Step Derivation)}
\addcontentsline{toc}{section}{Main Proof}
We now prove Theorem \ref{thm:global}.  Throughout the proof we denote 
$g_{k}:=\nabla f(x_{k})$ and $p_{k}$ the step produced by the algorithm.

\subsection*{Step 1 – Relating actual and predicted reduction}
From Lemma~\ref{lem:descent} and the definition of $m_{k}$,
\[
f(x_{k}+p_{k})\le f(x_{k})+g_{k}^{\top}p_{k}+\frac{L}{2}\|p_{k}\|^{2}
= m_{k}(p_{k})+\frac{L}{2}\|p_{k}\|^{2}.
\]
Hence
\[
\operatorname{ared}_{k}=f(x_{k})-f(x_{k}+p_{k})\ge
m_{k}(0)-m_{k}(p_{k})-\frac{L}{2}\|p_{k}\|^{2}
= \operatorname{pred}_{k}-\frac{L}{2}\|p_{k}\|^{2}.
\tag{1}
\label{eq:ared}
\]

\subsection*{Step 2 – Lower bound on the predicted reduction}
Because $p_{k}$ is at least a Cauchy point, Lemma~\ref{lem:cauchy} gives
\[
\operatorname{pred}_{k}=m_{k}(0)-m_{k}(p_{k})
\;\ge\;
\frac{1}{2}\,\frac{\|g_{k}\|^{2}}{\|B_{k}\|+\frac{\|g_{k}\|}{\Delta_{k}}\;\Delta_{k}.
\tag{2}
\label{eq:pred}
\]

\subsection*{Step 3 – Bounding $\|p_{k}\|$}
By construction $\|p_{k}\|\le\Delta_{k}$.  Using this together with \eqref{eq:ared} we obtain
\[
\operatorname{ared}_{k}\;\ge\;
\operatorname{pred}_{k}-\frac{L}{2}\Delta_{k}^{2}.
\tag{3}
\label{eq:ared2}
\]

\subsection*{Step 4 – Acceptance condition}
If $\rho_{k}\ge\eta_{1}$ the step is accepted, i.e. $x_{k+1}=x_{k}+p_{k}$.
From the definition of $\rho_{k}$ and \eqref{eq:ared2},
\[
\rho_{k}\ge\eta_{1}\;\Longrightarrow\;
\operatorname{pred}_{k}\le\frac{1}{\eta_{1}}\operatorname{ared}_{k}
\le\frac{1}{\eta_{1}}\bigl(\operatorname{pred}_{k}+\frac{L}{2}\Delta_{k}^{2}\bigr).
\]
Rearranging yields a bound on the predicted reduction for any accepted step:
\[
\operatorname{pred}_{k}\;\le\;\frac{L}{2}\,\frac{1}{1-\eta_{1}}\;\Delta_{k}^{2}.
\tag{4}
\label{eq:predbound}
\]

\subsection*{Step 5 – Sufficient decrease per successful iteration}
Combine \eqref{eq:pred} and \eqref{eq:predbound}.  For a successful iteration we have
\[
\frac{1}{2}\,\frac{\|g_{k}\|^{2}}{\|B_{k}\|+\frac{\|g_{k}\|}{\Delta_{k}}}\,\Delta_{k}
\;\le\;\operatorname{pred}_{k}
\;\le\;\frac{L}{2}\,\frac{1}{1-\eta_{1}}\;\Delta_{k}^{2}.
\]
Cancelling $\tfrac12\Delta_{k}$ and solving for $\|g_{k}\|$ gives
\[
\|g_{k}\|\;\le\;
\Bigl(\|B_{k}\|+\frac{\|g_{k}\|}{\Delta_{k}}\Bigr)
\frac{L}{1-\eta_{1}}\;\Delta_{k}.
\]
Since $\|B_{k}\|\le L$ (by $L$‑Lipschitzness of the gradient) and $\Delta_{k}\le\Delta_{\max}$,
\[
\|g_{k}\|\;\le\;
\Bigl(L+\frac{\|g_{k}\|}{\Delta_{\min}}\Bigr)\frac{L}{1-\eta_{1}}\;\Delta_{\max}.
\tag{5}
\label{eq:gradbound}
\]
Inequality \eqref{eq:gradbound} can be rearranged to obtain a *uniform* lower bound on the predicted reduction in terms of the gradient norm:
\[
\operatorname{pred}_{k}\;\ge\;
c_{1}\,\|g_{k}\|\,\Delta_{k},
\qquad
c_{1}:=\frac{1}{2}\,\frac{1}{\|B_{k}\|+\frac{\|g_{k}\|}{\Delta_{k}}}
\ge\frac{1}{2}\,\frac{1}{L+\frac{1}{\Delta_{\min}}\,\|g_{k}\|}.
\]
For all $k$ with $\|g_{k}\|>\varepsilon$ we have
\[
\operatorname{pred}_{k}\;\ge\;
c_{1}(\varepsilon)\,\varepsilon\,\Delta_{\min},
\quad\text{where }c_{1}(\varepsilon):=
\frac{1}{2}\,\frac{1}{L+\varepsilon/\Delta_{\min}}.
\tag{6}
\label{eq:pred_eps}
\]

\subsection*{Step 6 – Summation of decreases}
Let $\mathcal{S}$ be the set of successful iterations (those with $\rho_{k}\ge\eta_{1}$).  
For $k\in\mathcal{S}$,
\[
f(x_{k})-f(x_{k+1})=\operatorname{ared}_{k}\ge\eta_{1}\operatorname{pred}_{k}.
\tag{7}
\label{eq:decrease}
\]
Summing \eqref{eq:decrease} over all successful iterations up to $K$,
\[
\sum_{k\in\mathcal{S},\,k\le K}\operatorname{pred}_{k}
\;\le\;\frac{1}{\eta_{1}}\bigl(f(x_{0})-f_{\inf}\bigr).
\tag{8}
\label{eq:sum_pred}
\]

\subsection*{Step 7 – Bounding the number of ``large‑gradient'' successful steps}
Assume that during the first $N$ successful iterations the gradient norm stays above $\varepsilon$, i.e. $\|g_{k}\|>\varepsilon$ for those steps.  
Using \eqref{eq:pred_eps} in \eqref{eq:sum_pred} yields
\[
N\;c_{1}(\varepsilon)\,\varepsilon\,\Delta_{\min}
\;\le\;\frac{1}{\eta_{1}}\bigl(f(x_{0})-f_{\inf}\bigr).
\]
Thus
\[
N\;\le\;
\frac{f(x_{0})-f_{\inf}}{\eta_{1}\,c_{1}(\varepsilon)\,\varepsilon\,\Delta_{\min}}
=
\frac{2\bigl(f(x_{0})-f_{\inf}\bigr)}{\eta_{1}\,\Delta_{\min}}\;
\Bigl(L+\frac{\varepsilon}{\Delta_{\min}}\Bigr)\,
\frac{1}{\varepsilon}.
\]
Since $\varepsilon\le L\Delta_{\min}$ (otherwise the term $\varepsilon/\Delta_{\min}$ dominates),
the dominant term is $O\!\bigl(\frac{L\,(f(x_{0})-f_{\inf})}{\varepsilon^{2}}\bigr)$.
This proves the iteration‑complexity claim.

\subsection*{Step 8 – Limit inferior of the gradient norm}
If infinitely many successful iterations occur, the bound above forces
$\|g_{k}\|\to0$ along a subsequence; otherwise the algorithm would terminate after a finite number of successful steps because the radius would eventually be reduced below $\Delta_{\min}$ and the step would be rejected, contradicting the radius lower bound. Hence
\[
\liminf_{k\to\infty}\|g_{k}\|=0.
\]
\hfill$\square$

%%%%%%%%%%%%%%%%%%%%%%%%%%%%%%%%%%%%%%%%%%%%%%%%%%%%%%%%%%%%%%%%%%%%%%%%%%%%
\section*{Rate Derivation (With Constants)}
\addcontentsline{toc}{section}{Rate Derivation}
From Step 7 we have the explicit constant
\[
C:=\frac{2\bigl(f(x_{0})-f_{\inf}\bigr)}{\eta_{1}\,\Delta_{\min}}\,
\Bigl(L+\frac{1}{\Delta_{\min}}\Bigr).
\]
Consequently,
\[
\#\{k\le K:\|g_{k}>\varepsilon\}\;\le\;\frac{C}{\varepsilon^{2}}.
\]
Thus the method attains an $\varepsilon$‑first‑order stationary point in at most
$\mathcal{O}\!\bigl(\varepsilon^{-2}\bigr)$ successful iterations, and the total number of iterations is of the same order because the number of unsuccessful iterations is bounded by a constant multiple of the number of successful ones (standard trust‑region analysis).

%%%%%%%%%%%%%%%%%%%%%%%%%%%%%%%%%%%%%%%%%%%%%%%%%%%%%%%%%%%%%%%%%%%%%%%%%%%%
\section*{Special Cases}
\addcontentsline{toc}{section}{Special Cases}
\begin{enumerate}[noitemsep]
\item \textbf{Quadratic functions.} If $f(x)=\tfrac12 x^{\top}Qx+b^{\top}x+c$ with $Q$ symmetric, the model $m_{k}$ is exact ($B_{k}=Q$). The Cauchy point coincides with the exact solution of \eqref{eq:subprob} when $Q$ is positive semidefinite, leading to finite‑step convergence.
\item \textbf{Convex $L$‑smooth case.} The bound reduces to the classic $O(1/\varepsilon^{2})$ rate for gradient descent, but the trust‑region mechanism automatically selects a stepsize that never exceeds $1/L$, guaranteeing stability without line‑search.
\item \textbf{Large‑scale problems.} When $B_{k}$ is replaced by a cheap approximation (e.g., $B_{k}=0$), the Cauchy step becomes a scaled gradient step with stepsize $\Delta_{k}/\|g_{k}\|$. The same analysis carries through because Lemma~\ref{lem:cauchy} holds with $\|B_{k}\|=0$.
\end{enumerate}

%%%%%%%%%%%%%%%%%%%%%%%%%%%%%%%%%%%%%%%%%%%%%%%%%%%%%%%%%%%%%%%%%%%%%%%%%%%%
\section*{Discussion}
\addcontentsline{toc}{section}{Discussion}
The trust‑region gradient method presented above inherits the robustness of classical trust‑region schemes while requiring only first‑order information (the Cauchy point). Under the minimal smoothness assumptions \ref{as:smooth}–\ref{as:radius}, we proved global convergence to first‑order stationary points and derived the optimal $O(\varepsilon^{-2})$ complexity bound that matches the best known rates for first‑order methods on non‑convex $L$‑smooth functions. The explicit constants in the bound clarify how the hyper‑parameters $\eta_{1},\Delta_{\min},\Delta_{\max}$ influence the number of iterations, providing practical guidance for tuning.

\end{document}